\documentclass{article}
\usepackage{amsmath}
\begin{document}
\title{Rendering the formation cue distinguishes upstream visibility loss from downstream clipping}
\maketitle

\paragraph{What is executed.}
The study renders finite planar patches into a real floating-point
depth framebuffer using headless software OpenGL. It runs the two
unchanged pinned Ogre2 depth vertex/fragment shader pairs on that
buffer. An authored camera supplies a centered pinhole projection and
quad frustum rays: the pin uses the listed model's width, height,
horizontal field of view, and nominal camera near/far plane, plus the
inspected parent renderer's outer distances and projection formula.
No lens/noise, full Ogre compositor, ROS, native scene loading, hardware
or measured CoSTAR package runtime is claimed. The scene is designed
for diagnosis, not an independent cave or mission.

The left patch at 5 supplies the old image's known pixel. The
center patch is at 7 or 8 for the near site and at 37 or 38 for
the far site; the exact world selects the target depth. All
clearances/coordinates are in this authored optical model, not a
rasterized body path. Camera offsets are alternate fixed
extrinsics, not claimed reachable extra motions during the
one-dimensional formation.

\paragraph{Two physically different causes of no usable cue.}
An independent pinhole ray test checks every pixel of every frame,
including nearest-patch selection, absent returns, and the calibration
between point axes and depth. In the 32 cases over scene, camera and
renderer settings, 2,457,600 ray comparisons agree with the finite
mask oracle. For finite points the largest observed coordinate
error is under .010 in these fixtures, not a device precision bound.

At renderer far 300 the open far scenes give distinct raw
observations (X about 37.00957 and 37.99674 at the nominated ray).
At renderer far 10 the target is absent from the original
depth buffer and the two full rendered frames are equal.
Likewise a physical nearer panel occluding the entire target
makes the frames for the two worlds equal even when renderer
far is 300. Extra radio or a pre-clipping tap cannot recover
the hidden world in either case. This differs from software
depth clip 10, where the raw values survived and the old image
branch removed the distinction only later.

A narrow panel offers a separate clearance control. With one
tested camera position, the complete frame agrees byte-for-byte
with the panel-free scene, not just at the selected ray. Moving
the camera laterally by .009 in an otherwise identical setup
makes the nominated ray hit the nearer panel instead of the
far target. Seeing no return at that location is not a proof
of absence of the panel; a fresh camera may also lose the
particular target while detecting something useful for another task.
No guarantee about arbitrary subpixel coverage follows.

\paragraph{Same-chain follow-through.}
Pass actual rendered RGBA values for the two named pixels into
the selected native PostRender callback flow, crop width 2, preserving
the identical shader values. The native depth emission and stock
point masking and pack then run under the chosen held subscription.
The original point label uses the decoded coordinate; the alternative
authorized source read interprets a narrow interval rather than
requiring exact integer arithmetic on graphics values. The read is
scope-bound to the stationary capture. Other staleness and membership
controls from the parent are not weakened by this change.
Only a valid bit consumes radio capacity, and the earlier radio
obligation and old source depth emission count separately.
Keep the same exact deadline model, first eligible send calendar and
live receive budget. The original source comparator sees both stock
outputs; the direct corrected comparator has equally qualified raw
access, and ties.

All 768 joins retain the actual old broker callback. With old depth
emission required, 40 scheduled cases pass with the original data path
and 64 with the scoped read. For the open far site at render far 300,
the earlier old-report deadline is still necessary at receive cap 2;
at cap 3 the slot exists, whereas closing visibility or setting render
far 10 defeats either budget. The relative availability of these modes
is not a measured frequency. The exact four motion proofs preserve
their declared requirements; no part of the scene test is presented
as an independently selected mission, deployed extension or hardware
clearance guarantee.
\end{document}
